\documentclass[11pt,a4paper]{article}

% ---------- packages (all standard in a normal TeX Live / MiKTeX setup) ----------
\usepackage[T1]{fontenc}
\usepackage{mathptmx}              % Times New Roman look
\usepackage[margin=2cm]{geometry}
\usepackage{booktabs}
\usepackage{longtable}
\usepackage{tabularx}
\usepackage{array}
\usepackage{enumitem}
\usepackage{amsmath,amssymb}
\usepackage{verbatim}
\usepackage{float}
\usepackage[hidelinks]{hyperref}

\setlength{\parskip}{4pt}
\emergencystretch=1.5em
\renewcommand{\arraystretch}{1.12}
\newcommand{\code}[1]{\texttt{#1}}

\begin{document}

% =====================================================================
% Rubric sheet (serves as the title page, as required by the handout)
% =====================================================================
\noindent
{\Large\bfseries DEPARTMENT OF COMPUTER AND INFORMATION SYSTEMS ENGINEERING}\\[2mm]
{\large BACHELORS IN COMPUTER SYSTEMS ENGINEERING}\\[1mm]
{\large Course Code and Title: CS-323 Artificial Intelligence}\\[1mm]
{\large Complex Engineering Problem --- Term Project}\\[1mm]
{\small TE Batch 2024, Fall Semester 2026 \quad Instructors: Dr.\ Maria Waqas, Dr.\ Sumayya Zafar}

\vspace{4mm}
\noindent
\setlength{\tabcolsep}{1.5pt}
\begin{tabular}{|p{1.2cm}|p{2.2cm}|p{1.4cm}|p{2.2cm}|p{1.4cm}|p{2.2cm}|p{1.4cm}|p{2.2cm}|p{1.4cm}|}
\hline
\bfseries Group & \multicolumn{2}{c|}{\bfseries S1}
                & \multicolumn{2}{c|}{\bfseries S2}
                & \multicolumn{2}{c|}{\bfseries S3}
                & \multicolumn{2}{c|}{\bfseries S4} \\
\cline{2-9}
                & \bfseries Name & \bfseries Roll
                & \bfseries Name & \bfseries Roll
                & \bfseries Name & \bfseries Roll
                & \bfseries Name & \bfseries Roll \\
\cline{2-9}
\footnotesize Members & Ukasha & CS-24115
        & Abdul Mueed & CS-24117
        & Abdullah Baig & CS-24123
        & Umer Iqbal & CS-24127 \\
\hline
\end{tabular}

\vspace{4mm}
\noindent{\bfseries Grading Rubric}\\[1mm]
\footnotesize
\setlength{\tabcolsep}{2pt}
\begin{longtable}{|p{3.3cm}|p{2.0cm}|p{2.0cm}|p{2.0cm}|p{2.0cm}|p{0.8cm}|p{0.8cm}|p{0.8cm}|p{0.8cm}|}
\hline
\bfseries Criterion and scale & \bfseries 1 & \bfseries 2 & \bfseries 3 & \bfseries 4
& \bfseries S1 & \bfseries S2 & \bfseries S3 & \bfseries S4 \\
\hline
\endhead
C1. Graph search (4 marks) &
Incorrect; wrong output or fails on boards &
Partial; limited or inconsistent results &
Correct; mostly correct outputs &
Correct and efficient &
& & & \\ \hline
C2. CSP formulation (4 marks) &
Incorrect model or solver &
Partly correct; limited or inaccurate output &
Correct model; mostly accurate results &
Accurate and efficient solver &
& & & \\ \hline
C3. Integration (2 marks) &
Not integrated; robot fails task &
Partial; misses visits or return &
Executes most valid cases &
Complete journey; all cleaned; returns &
& & & \\ \hline
C4. Comparison and analysis (2 marks) &
No meaningful comparison &
Basic comparison; little depth &
Clear metrics and interpretation &
Comprehensive quantitative/qualitative analysis &
& & & \\ \hline
C5. Report format (4 marks) &
Missing information; poorly formatted &
Partial information; well formatted &
All information; poorly formatted &
All information; correct format &
& & & \\ \hline
C6. Team contribution (4 marks) &
No meaningful contribution &
Partial contribution; limited collaboration &
Adequate contribution; met goals &
Active, significant contribution beyond goals &
& & & \\ \hline
\bfseries Final score & \multicolumn{8}{c|}{
\hspace{1cm} Teacher's signature: \underline{\hspace{5cm}} } \\
\hline
\end{longtable}
\normalsize

\newpage

% =====================================================================
\section{Problem Definition and Representation}

The robot works on a grid of square cells. It starts at a charging station, must visit every
dirty cell exactly once as decided by a scheduler, and must finally return to the station.
Four-direction movement only (up, down, left, right); diagonal moves are not allowed.
Obstacle cells cannot be entered. Every movement costs one unit, so the navigation cost of a
journey is simply the number of moves.

The two artificial-intelligence components are:

\begin{itemize}[leftmargin=1.2cm]
  \item \textbf{Graph-based search}: find a valid shortest path from a given start cell to a
        given goal cell while avoiding walls. We implement one uninformed method (BFS) and one
        informed method (A*).
  \item \textbf{Constraint Satisfaction Problem (CSP)}: assign each cleaning task to one time
        slot so that all precedence, latest-slot and one-task-per-slot rules hold.
\end{itemize}

The CSP output fixes the \emph{order} of cleaning; the search component supplies the
\emph{paths}. The two are then integrated into one simulation.

\subsection{Problem formulation}

\begin{center}
\begin{tabular}{@{}p{3.4cm}p{12.3cm}@{}}
\toprule
Element & Representation \\
\midrule
Environment & Dictionary with \code{rows}, \code{cols}, \code{obstacles} (set of cells),
              \code{dirties} (ordered list) and \code{station}, all read from \code{Environment.txt}. \\
State & A cell $(r,c)$, with $0\le r<\text{rows}$ and $0\le c<\text{cols}$. \\
Initial state & The cell passed to the search function (the charging station in the full run). \\
Actions & UP $(-1,0)$, DOWN $(+1,0)$, LEFT $(0,-1)$, RIGHT $(0,+1)$. \\
Transition model & New state $= (r+\Delta r,\; c+\Delta c)$; valid only inside the grid and not a wall. \\
Goal test & \code{state == goal\_state}. \\
Step cost & Every valid move costs 1; therefore path cost $=\;|\text{path}|-1$. \\
CSP variables & One per dirty cell: $D1,D2,\dots,Dn$, in file order. \\
CSP domains & Initially $\{1,2,\dots,n_{\text{slots}}\}$; reduced by latest-slot rules. \\
CSP constraints & Precedence ($D_i$ before $D_j$), latest slot ($D_i\le T_k$),
                  and all-different slot assignment. \\
\bottomrule
\end{tabular}
\end{center}

\subsection{Example environment (our \code{Environment.txt})}

The 8$\times$8 grid has seven walls, four dirty cells and one charging station at $(0,0)$.
Letters $D1\ldots D4$ are the dirty cells, $X$ is a wall, $C$ the station and ``.'' a free cell.

\begin{center}
\small
\begin{tabular}{c|cccccccc}
      & 0 & 1 & 2 & 3 & 4 & 5 & 6 & 7 \\ \hline
0 & C & . & . & . & . & . & . & . \\
1 & . & . & X & . & . & . & . & D1 \\
2 & . & . & X & . & . & . & . & . \\
3 & . & . & X & . & . & . & . & . \\
4 & . & . & X & . & . & X & . & D2 \\
5 & . & . & X & . & . & X & . & . \\
6 & D4& . & . & . & . & . & . & . \\
7 & . & . & . & . & . & . & D3& . \\
\end{tabular}
\end{center}

% =====================================================================
\section{Graph-Based Search Algorithms}

\subsection{Breadth-First Search (uninformed)}

BFS keeps a FIFO queue of cells to explore. It expands all cells at depth $d$ before any cell
at depth $d+1$. Because every step has cost 1, the first time a cell is generated it is reached
by a shortest path; a parent map records that path. BFS is complete on a finite grid and
optimal for unit costs, but it may expand a large fraction of the grid before reaching the goal.

\begin{verbatim}
BFS(grid, start, goal):
    queue <- [start];  parent[start] <- None;  expanded <- 0
    while queue is not empty:
        cell <- queue.pop_front();  expanded <- expanded + 1
        if cell = goal: return build_path(parent, goal), cost, expanded
        for each legal neighbour nxt of cell:
            if nxt not in parent:
                parent[nxt] <- cell
                queue.push_back(nxt)
    return FAIL, FAIL, expanded
\end{verbatim}

\subsection{A* Search (informed)}

A* orders the frontier by $f(n)=g(n)+h(n)$, where $g(n)$ is the cost already paid to reach $n$
and $h(n)$ is the Manhattan distance to the goal:
\[
h(n)=|r_n-r_{\text{goal}}|+|c_n-c_{\text{goal}}|.
\]
Manhattan distance is \emph{admissible} (it never overestimates the true remaining cost in a
four-direction grid with unit steps) and \emph{consistent} (moving one step changes $h$ by at
most 1, which is exactly the step cost). Therefore the first time A* pops a goal cell, that
goal is optimal, and no already-expanded cell needs to be reopened. With a good estimate the
search is strongly directed toward the target and expands far fewer cells than BFS.

\begin{verbatim}
A_STAR(grid, start, goal):
    todo <- [(h(start), start)];  parent[start] <- None
    cost[start] <- 0;  done <- {};  expanded <- 0
    while todo is not empty:
        sort todo by estimated cost;  (f, cell) <- todo.pop_front()
        if cell in done: continue
        done.add(cell);  expanded <- expanded + 1
        if cell = goal: return build_path(parent, goal), cost, expanded
        for each legal neighbour nxt of cell:
            new_cost <- cost[cell] + 1
            if nxt not in cost or new_cost < cost[nxt]:
                cost[nxt] <- new_cost;  parent[nxt] <- cell
                todo.push((new_cost + h(nxt), nxt))
    return FAIL, FAIL, expanded
\end{verbatim}

(For these small grids we sort the frontier list explicitly; this is slower in theory than a
binary heap but much easier to read, and the grids are at most a few hundred cells.)

\subsection{Property comparison}

\begin{center}
\begin{tabular}{@{}p{3.0cm}p{5.8cm}p{5.8cm}@{}}
\toprule
Property & BFS & A* \\
\midrule
Type & Uninformed & Informed (uses Manhattan $h$) \\
Frontier & FIFO queue & List sorted by $f=g+h$ \\
Complete & Yes (finite grid) & Yes (finite grid) \\
Optimal & Yes, unit step costs & Yes, admissible and consistent $h$ \\
Typical expansions & High (layer by layer) & Low (goal-directed) \\
Memory & $O(b^d)$ frontier & $O(b^d)$ frontier \\
\bottomrule
\end{tabular}
\end{center}

% =====================================================================
\section{CSP Formulation and Implementation}

\subsection{Variables, domains and constraints}

\begin{itemize}[leftmargin=1.2cm]
  \item \textbf{Variables}: $D1,D2,\ldots,Dn$ --- one for each dirty location in
        \code{Environment.txt}. A task that appears in no scheduling rule is still a variable.
  \item \textbf{Domains}: integers $1\ldots N$ ($N$ = number of time slots); a latest rule
        $D_i$ by $T_k$ trims the domain of $D_i$ to $\{1\ldots k\}$.
  \item \textbf{Precedence}: rule $(D_i,D_j)$ means slot$(D_i) <$ slot$(D_j)$.
  \item \textbf{Latest slot}: rule $(D_i,T_k)$ means slot$(D_i)\le k$.
  \item \textbf{Slot exclusivity}: all assigned slots are distinct, so $N\ge n$ is required.
\end{itemize}

\subsection{Role of each required function}

\begin{center}
\footnotesize
\begin{tabular}{@{}p{4.6cm}p{11.1cm}@{}}
\toprule
Function & Role in the solver \\
\midrule
\code{select\_unassigned\_variable} & MRV: choose the unassigned variable with the smallest
current domain. Ties are broken by the degree heuristic: the variable linked by precedence
rules to the largest number of other unassigned variables wins; a final alphabetical tie-break
keeps runs deterministic. \\
\code{order\_domain\_values} & Least Constraining Value: count how many values of related
unassigned variables would be removed (by slot exclusivity or by precedence) if this value
were chosen, and try the least damaging value first. \\
\code{is\_consistent} & Reject a value if the slot is already used, if it violates a latest-slot
rule, or if it breaks a precedence rule whose partner is already assigned. \\
\code{forward\_check} & Work on a fresh copy of the domains, remove the chosen slot from other
variables, enforce precedence bounds on directly linked variables, and reject the branch if any
domain becomes empty. Fresh copies give automatic restoration on backtracking. \\
\code{backtracking\_search} & Recursive depth-first assignment using the four helpers above;
returns a complete assignment or failure. \\
\bottomrule
\end{tabular}
\end{center}

\begin{verbatim}
BACKTRACKING(variables, domains, rules):
    assignment <- {}
    return SOLVE(domains)

SOLVE(domains):
    if all variables assigned: return copy of assignment
    var  <- SELECT_UNASSIGNED_VARIABLE(assignment, domains, rules)   % MRV + degree
    vals <- ORDER_DOMAIN_VALUES(var, assignment, domains, rules)     % LCV
    for value in vals:
        if not IS_CONSISTENT(var, value, assignment, rules): continue
        smaller <- FORWARD_CHECK(var, value, assignment, domains, rules)
        if smaller = FAIL: continue
        assignment[var] <- value
        answer <- SOLVE(smaller)
        if answer != FAIL: return answer
        delete assignment[var]
    return FAIL
\end{verbatim}

\subsection{Sample schedule produced}

For our \code{Schedule.txt} ($D1\!\prec\!D2$, $D2\!\prec\!D3$, $D1\!\prec\!D4$;
$D2\le T3$, $D3\le T4$) the solver returns:

\begin{center}
\begin{tabular}{c|cccc}
Time slot & T1 & T2 & T3 & T4 \\ \hline
Task      & D1 & D2 & D3 & D4 \\
Dirty cell& (1,7) & (4,7) & (7,6) & (6,0) \\
\end{tabular}
\end{center}

All four tasks are assigned once, all slots are distinct, and every precedence/latest rule is
satisfied. Notice that the sample files yield a different valid order
($D2,D1,D3,D4$); the solver is not hard-coded to one answer.

% =====================================================================
\section{System Integration}

\code{execute\_schedule(grid, charging\_station, schedule, algorithm)} performs the full
simulation:

\begin{enumerate}[leftmargin=1.2cm]
  \item Sort the CSP assignment by slot number to obtain the cleaning order.
  \item Translate each task label into its dirty-cell coordinate using the environment file.
  \item Starting at the charging station, call \code{search} for every consecutive stop.
  \item Append each returned path while skipping its first cell (already present), so the
        complete route has no duplicated joints.
  \item Add the corresponding move cost of every leg.
  \item After the last dirty cell, call \code{search} once more back to the charging station.
  \item Verify that all dirty cells appear in the schedule and that the final cell is the station.
        Return \code{complete\_path}, \code{total\_cost} and \code{success}.
  \item If any leg returns failure, stop immediately and report failure rather than claiming
        that cleaning succeeded.
\end{enumerate}

The components remain independent: the CSP solver never sees the grid, and the search
functions never see the time slots. Integration only consumes their outputs.

% =====================================================================
\section{Results and Comparative Analysis}

\subsection{Test 1 --- our own environment}

\begin{center}
\begin{tabular}{@{}p{3.2cm}rrrr@{}}
\toprule
Journey & BFS moves & BFS expanded & A* moves & A* expanded \\
\midrule
(0,0) $\to$ (1,7) & 8 & 38 & 8 & 15 \\
(1,7) $\to$ (4,7) & 3 & 10 & 3 & 4 \\
(4,7) $\to$ (7,6) & 4 & 18 & 4 & 7 \\
(7,6) $\to$ (6,0) & 7 & 38 & 7 & 8 \\
(6,0) $\to$ (0,0) & 6 & 24 & 6 & 7 \\
\midrule
\textbf{Total} & \textbf{28} & \textbf{128} & \textbf{28} & \textbf{41} \\
\bottomrule
\end{tabular}
\end{center}

Complete A* route (29 cells, 28 moves):
\begin{verbatim}
(0,0) (0,1) (0,2) (0,3) (0,4) (0,5) (0,6) (0,7) (1,7)
(2,7) (3,7) (4,7) (4,6) (5,6) (6,6) (7,6) (6,6) (6,5)
(6,4) (6,3) (6,2) (6,1) (6,0) (5,0) (4,0) (3,0) (2,0)
(1,0) (0,0)
\end{verbatim}

\subsection{Test 2 --- provided sample environment}

\begin{center}
\begin{tabular}{@{}p{3.2cm}rrrr@{}}
\toprule
Journey & BFS moves & BFS expanded & A* moves & A* expanded \\
\midrule
(0,0) $\to$ (3,7) & 10 & 48 & 10 & 29 \\
(3,7) $\to$ (0,7) & 3 & 10 & 3 & 4 \\
(0,7) $\to$ (7,7) & 7 & 27 & 7 & 8 \\
(7,7) $\to$ (7,0) & 7 & 33 & 7 & 8 \\
(7,0) $\to$ (0,0) & 7 & 26 & 7 & 8 \\
\midrule
\textbf{Total} & \textbf{34} & \textbf{144} & \textbf{34} & \textbf{57} \\
\bottomrule
\end{tabular}
\end{center}

Complete A* route for the sample run (35 cells, 34 moves):
\begin{verbatim}
(0,0) (0,1) (0,2) (0,3) (0,4) (0,5) (0,6) (0,7) (1,7)
(2,7) (3,7) (2,7) (1,7) (0,7) (1,7) (2,7) (3,7) (4,7)
(5,7) (6,7) (7,7) (7,6) (7,5) (7,4) (7,3) (7,2) (7,1)
(7,0) (6,0) (5,0) (4,0) (3,0) (2,0) (1,0) (0,0)
\end{verbatim}

\subsection{Interpretation}

\begin{itemize}[leftmargin=1.2cm]
  \item On every leg both algorithms return the \emph{same optimal cost}; this is exactly what
        theory predicts for unit step costs and an admissible, consistent heuristic.
  \item A* expands roughly one third of the cells BFS expands (41 vs 128, and 57 vs 144 in the
        two full journeys), because Manhattan distance guides it toward the target instead of
        filling complete layers.
  \item On very short or unobstructed legs the gap is small; it grows with the distance and
        with obstacles that force detours.
  \item The two total costs (28 and 34) differ because they belong to different environments
        and different CSP-generated orders, not because one algorithm is ``better'' at finding
        short routes. The comparison that matters is per-leg cost (equal) and expansions
        (A* lower).
\end{itemize}

% =====================================================================
\section{Testing and Validation}

Three categories of test were run from the notebook (additional edge cases: start equals goal,
unreachable target, hand-written rules containing \code{'T3'} strings, and coordinates wrapped
over several lines).

\subsection{Test 1 --- own environment and schedule}

Configuration:
\begin{verbatim}
Environment.txt                                 Schedule.txt
8,8                                             4
7                                               3
(1,2),(2,2),(3,2),(4,2),(5,2),(4,5),(5,5)       (D1,D2)
4                                               (D2,D3)
(1,7),(4,7),(7,6),(6,0)                         (D1,D4)
(0,0)                                           2
                                                (D2,T3)
                                                (D3,T4)
\end{verbatim}

Observed output: schedule T1=D1, T2=D2, T3=D3, T4=D4; order
$(1,7)\!\to\!(4,7)\!\to\!(7,6)\!\to\!(6,0)$; complete route as listed above;
28 movements, total navigation cost 28; status SUCCESS --- all dirty cells cleaned and the
robot returned to $(0,0)$. The notebook displays a coloured grid of the final state (all four
dirty cells marked cleaned, robot at C).

\begin{figure}[H]
\centering
\fbox{\parbox{0.92\textwidth}{\small
\textbf{Screenshot placeholder --- Test 1.} Paste here the notebook screenshot showing
\code{Environment.txt}/\code{Schedule.txt} and the SUCCESS output with the coloured journey grid.
}}
\caption{Configuration and successful output for our environment (notebook output).}
\end{figure}

\subsection{Test 2 --- provided sample files}

Configuration: \code{Environment\_Sample.txt} (8$\times$8, six walls in column 3,
dirty cells $(0,7),(3,7),(7,7),(7,0)$, station $(0,0)$) and
\code{Schedule\_Sample.txt} (four slots; $D1\!\prec\!D3$, $D2\!\prec\!D4$, $D1\!\prec\!D4$;
$D3\le T3$, $D2\le T3$).

Observed output: schedule T1=D2, T2=D1, T3=D3, T4=D4 (also valid); order
$(3,7)\!\to\!(0,7)\!\to\!(7,7)\!\to\!(7,0)$; 34 movements, total cost 34; status SUCCESS.

\begin{figure}[H]
\centering
\fbox{\parbox{0.92\textwidth}{\small
\textbf{Screenshot placeholder --- Test 2.} Paste here the notebook screenshot of the sample
run with its comparison chart and final grid.
}}
\caption{Configuration and successful output for the provided sample (notebook output).}
\end{figure}

\subsection{Test 3 --- failure cases}

\noindent\textbf{3a. Infeasible schedule} --- \code{Environment.txt} with
\code{Schedule\_Error.txt}: three slots for four tasks, plus $D1\!\prec\!D2$, $D2\!\prec\!D3$,
$D2\le T2$, $D3\le T2$. The precedence chain forces $D2\ge 2$ and $D3\ge 3$, which contradicts
$D3\le 2$. The CSP correctly reports \emph{No feasible schedule exists}; no navigation is
attempted.

\noindent\textbf{3b. Unreachable destination} --- \code{Environment\_Error.txt} with
\code{Schedule.txt}: dirty cell $D4=(6,1)$ is enclosed by walls at
$(5,0),(5,1),(5,2),(6,0),(6,2),(7,0),(7,1),(7,2)$. The CSP still finds a schedule, but both
BFS and A* return no path for $(7,7)\to(6,1)$ after expanding 55 cells. The simulation stops
with a partial route of 14 moves and status \emph{FAIL --- a required path was not found};
it does not falsely report success.

Final grid of the failure run ($R$ = robot stuck at $(7,7)$, $*$ = cleaned, $X$ = wall):
\begin{center}
\small
\begin{tabular}{c|cccccccc}
      & 0 & 1 & 2 & 3 & 4 & 5 & 6 & 7 \\ \hline
0 & C & . & . & . & . & . & . & * \\
1 & . & . & . & . & . & . & . & . \\
2 & . & . & . & . & . & . & . & . \\
3 & . & . & . & . & . & . & . & * \\
4 & . & . & . & . & . & . & . & . \\
5 & X & X & X & . & . & . & . & . \\
6 & X & * & X & . & . & . & . & . \\
7 & X & X & X & . & . & . & . & R \\
\end{tabular}
\end{center}

\begin{figure}[H]
\centering
\fbox{\parbox{0.92\textwidth}{\small
\textbf{Screenshot placeholder --- Test 3.} Paste here notebook screenshots showing the
infeasible-schedule message and the unreachable-destination FAIL status with its grid.
}}
\caption{Detection and reporting of both failure cases.}
\end{figure}

\subsection{Test 4 --- infeasible input}

The input files and query points are checked before a run starts, and an infeasible case is
reported on the console instead of crashing or producing a false success:

\begin{itemize}[leftmargin=1.2cm]
  \item A bad environment or schedule file (missing, short, with extra numbers, an invalid grid
        size, overlapping or out-of-grid cells, an unknown task, etc.) prints
        \code{INPUT ERROR: <reason>} and the run stops.
  \item A schedule with too few slots, an impossible deadline or contradictory rules prints
        \code{NO FEASIBLE SCHEDULE} (as in Test 3a above).
  \item A walled-in target, a start/goal on a wall or outside the grid, or an unknown algorithm
        prints \code{no path} / \code{FAILED}, and the robot stops at the last reachable cell
        (as in Test 3b above).
\end{itemize}

The algorithms themselves (BFS, A*, MRV/LCV/forward checking and backtracking) are unchanged;
these are simple checks at the file-reading and \code{search} entry points. Normal edge cases
remain normal: \code{start = goal} returns \code{([start], 0, 0)}, an empty route has cost
\code{None}, two deadlines keep the stricter one, and unused slots stay empty.

% =====================================================================
\section{Discussion and Improvements}

\paragraph{Correctness.} Both searches return legal, adjacent, wall-free paths; the cost equals
path length minus one. Start equals goal is handled without expanding a node. When no route
exists, the algorithms return failure and the integration reports failure. The CSP assigns
every task exactly once and enforces all three constraint families.

\paragraph{Completeness and optimality.} BFS is complete and optimal for unit step costs on a
finite grid. A* is complete here and optimal because Manhattan distance is admissible and
consistent on a four-direction grid. Backtracking with MRV and forward checking is complete for
the finite CSP: it explores all consistent assignments and returns a solution if one exists.

\paragraph{Limitations.} Obstacles and tasks are static; step costs are all 1; battery level,
recharging needs, turning cost and dynamic obstacles are not modelled. The CSP optimises only
feasibility, not travel distance --- a feasible order may be longer than another feasible order.
The A* frontier is a sorted list rather than a heap, which is fine for small grids but does not
scale to very large maps.

\paragraph{Possible improvements.} Use a binary heap or bidirectional A*; use branch-and-bound
or a distance-aware objective so the CSP minimises total travel; add battery constraints and
opportunity charging; animate the journey; add unit tests for parsers and failure handling.

% =====================================================================
\section{References}

\begin{enumerate}[leftmargin=1.2cm]
  \item Hand Outs of Dr. Maria Waqas for CSP and BFS algorithms
  \item Python Software Foundation, \emph{Python 3 Documentation} --- \code{heapq},
        \code{deque}, \code{re}, \code{IPython.display}.
\end{enumerate}

% =====================================================================
\section{Generative AI Use Declaration}

Generative AI (the OpenCode coding assistant) was used during this project to help debug the
implementation, simplify the code for readability, generate the coloured notebook visuals, and
draft parts of this report. The group members reviewed the code, checked every result against
the handout requirements, and understand the algorithms and outputs. We take full
responsibility for the submitted work.

\end{document}
